\scanpage{43}
\mspage{- IV - \inserted{977} - \inserted{12} -}
$u:B\longrightarrow A$, où $B=k\inserted{[}[T_1,\ldots,T_r]\inserted{]}$ ($r$ entier quelconque) faisant de $A$ une $B$-algèbre finie, $\uline{m}$ est le seul point de $\operatorname{Spec}(A)$ au-dessus\unsure{'} \unsure{de|du} l'idéal maximal \deleted{\unsure{son}} $\uline{n}$ de $B$ ; \deleted{\unsure{donc}} le morphisme \deleted{$\unsure{{}^a u}:\operatorname{Spec}(A)\longrightarrow\operatorname{Spec}(B)$}\inserted{$\operatorname{Spec}(A)\longrightarrow\operatorname{Spec}(B)$} ne peut être non ramifié en $\uline{m}$ que si $A/\uline{n}A$ est un corps, extension séparable de $k$ (17.4.1), ce qui implique que le corps résiduel $K$ de $A$ est une extension séparable de $k$. Si cette condition n'est pas remplie, la conclusion de (18.11.5) ne peut jamais être vérifiée pour $\uline{p}=\uline{m}$, quel que soit l'entier $m$.

(ii) Supposons toujours $\uline{p}\neq\uline{m}$, mais considérons le cas où la condition de (18.11.5) relative à \inserted{l'existence de} l'idéal minimal $\uline{q}$ n'est pas remplie. Soit $\uline{J}$ le noyau de l'homomorphisme canonique $A\longrightarrow A_{\uline{p}}$, autrement dit l'ensemble des annulateurs des éléments de $A-\uline{p}$, et posons $A'=A/\uline{J}$ ; on a nécessairement $\uline{J}\subset\uline{p}$, et si l'on pose $\uline{p}'=\uline{p}/\uline{J}$, l'homomorphisme canonique $A_{\uline{p}}\longrightarrow A'_{\uline{p}'}$ est \uline{bijectif}. En vertu de (18.11.3.\unsure{3}) et de (\uline{0},20.4.12), l'homomorphisme canonique
\deleted{$\unsure{\widehat\Omega^1_{A/k}\otimes_A A'\longrightarrow\widehat\Omega^1_{A'/k}}$}
$\overset{\inserted{\scriptscriptstyle\wedge}}{\Omega}^1_{A/k}\otimes_A A'\longrightarrow\widehat\Omega^1_{A'/k}$ est \uline{surjectif}, donc il en est de même de $(\widehat\Omega^1_{A/k})_{\uline{p}}\longrightarrow(\widehat\Omega^1_{A'/k})_{\uline{p}'}$. D'autre part, l'homomorphisme canonique $A'\longrightarrow A'_{\uline{p}'}$ est par définition \uline{injectif}, et les idéaux \inserted{premiers} minimaux de $A'$ correspondent donc biunivoquement à ceux de $A'_{\uline{p}'}$ ; \deleted{\unsure{donc}}\inserted{par suite} Il y a un idéal minimal $\uline{q}'$ de $A'$ \deleted{\unsure{et|qui}} contenu dans \deleted{\unsure{un}} $\uline{p}'$ et tel que $\dim(A'/\uline{q}')=\dim(A')$. On peut ainsi appliquer (18.11.5) à $A'$ et $\uline{p}'$, et il y a donc un $k$-homomorphisme local $u':B\longrightarrow A'$, avec $B=k\inserted{[}[T_1,\ldots,T_m]\inserted{]}$ tel que le morphisme correspondant \deleted{$\unsure{{}^a u'}:$}\inserted{$f':$}$\operatorname{Spec}(A')\longrightarrow\operatorname{Spec}(B)$ soit fini et non ramifié au point $\uline{p}'$. Mais on peut factoriser $u'$ en $B\xrightarrow{u}A\xrightarrow{p}A'$ où \deleted{\unsure{$u$ est injectif et}} $p$ est l'homomorphisme canonique ; d'autre part, l'injection canonique $\operatorname{Spec}(A')\longrightarrow\operatorname{Spec}(A)$ est un\deleted{e} \deleted{\uline{\unsure{immersion ouverte|immersion fermée}}} \uline{isomorphisme local} au point $\uline{p}'$ en vertu de l'isomorphisme $A_{\uline{p}}\xrightarrow{\sim}A'_{\uline{p}'}$ et de (\uline{I},6.5.4). On voit donc que l'hypothèse a) de (18.11.5) entraîne encore, dans ce cas, l'existence d'un $k$-homomorphisme local $u:B\longrightarrow A$ tel que

\margin{\uline{\unsure{Virer|Vider}} (ii) sauf \unsure{raison} sérieuse !}
\mspage{42}
